\chapter{贪污贿赂罪}

% Source PDF physical page 397; see audit/source_coverage.json.
\begin{EditionOpeningPanel}
超级重点罪名：贪污罪、挪用公款罪、受贿罪。复习这三大罪名，方法是深刻理解其行为公式。
\end{EditionOpeningPanel}

\section{贪污罪}

第382条国家工作人员\CJKunderwave{利用职务上的便利}，侵吞、窃取、骗取或者以其他手段非法占有公共财物的，是贪污罪。

受国家机关、国有公司、企业、事业单位、人民团体委托管理、经营国有财产的人员，利用职务上的便利，侵吞、窃取、骗取或者以其他手段非法占有国有财物的，以贪污论。

与前两款所列人员勾结，伙同贪污的，以共犯论处。

贪污罪，是指国家工作人员利用职务上的便利，侵吞、窃取、骗取或者以其他手段非法占有公共财物的行为。贪污罪的保护法益有两个，一是公家的财产权，二是职务行为的廉洁性。

\subsection*{一、构成要件}

\InsertDiagram{figures/manual_table_p397_1.tex}

\subsubsection*{（一）行为主体}

行为主体是国家工作人员。国家工作人员是指依法从事公务的人员。

[总结]公司企业中，以下四类人属于国家工作人员：

1.国有企业中从事公务的人员。（第93条）

2.国有企业委派到非国有单位从事公务的人员。（第93条）

% Source PDF physical page 398; see audit/source_coverage.json.
3.国有控股公司、参股公司中哪些人是国家工作人员？经国有公司提名或批准，代表国有公司在国有控股公司、参股公司中从事管理工作的人员，是国家工作人员。（司法解释）

4.受国家单位委托，管理、经营国有财产的人员。这类人本不属于国家工作人员，但是，刑法第382条第2款将他们拟制为国家工作人员。由于挪用公款罪中没有类似法律拟制，因此此类人不是挪用公款罪中的“国家工作人员”，但可构成挪用资金罪。

[提示】通过伪造国家机关公文、证件担任了国家工作人员，可以成为贪污贿赂犯罪的主体。

\subsubsection*{（二）行为方式}

行为方式上必须利用职务上的便利。

1.这里的“职务”要求具有一定的管理性，不包括纯粹的体力劳动。

2.这里的“利用”要求实质利用，不包括形式利用。实质利用，是指利用职权发挥了实质影响力。形式利用，是指因工作关系熟悉作案环境，因工作关系易于接近作案目标，因工作关系容易进入某些场所。（2002年第9题）

\subsubsection*{（三）实行行为}

1.侵吞

这里的侵吞与侵占罪中的“侵占”含义相同，也即，将他人所有、自己占有的财物变成自己所有，在这里就是指，将单位所有、自己依职权占有的财物变成自己所有。

（1）占有的种类

第一，直接占有。例1，国家工作人员甲一个人带着公款出公差，将公款打入自已账户。甲构成贪污。例2，动植物检疫人员甲将自已负责检验的肉品，私自割下一块，拿回家让老婆烧菜。甲依据职权直接占有财物，变成自己所有，构成贪污。

第二，领导的松弛占有。这是指，领导没有直接占有财物，但对财物享有支配权，或对直接管理财物的职员享有管理权。例如，国企总经理甲指示出纳乙将国企一笔公款拿到甲家。甲属于将公家所有、自已依据职权松弛占有的财物变成自已所有，构成贪污罪。乙若知道实情，则构成贪污罪的共同正犯。\footnote{提示：单位对财物可以有所有权，但单位无法占有财物，因为“占有”要求有占有意思，所以占有人只能是自然人。例如，单位仓库里的财物属于单位所有，由库管员占有。又如，单位公账里的资金是单位所有，但不是单位占有。该资金的占有人包括两类人。一是直接占有人，也即负责管理该账户的会计、出纳等。二是松弛占有人，也即对该账户资金有支配权或对会计、出纳有管理权的领导，如财务总监、总经理等。这些人将公账资金转人自已账户，属于侵占行为，而非盗窃行为。}

（2）财物的种类

常见的公共财物是有体物。例如，国家工作人员在对外交往中接受礼物，应当交公而不交公，构成贪污。

[总结]钱款。第一，已收账款。例如，国企会计甲代表单位从客户处收到账款，然后不入公账，据为已有，构成侵占型贪污罪。

第二，应收账款。例如，乙欠某国企100万元，该国企总经理甲对乙讲：“我来收款，你将欠款的一半打给我，我将账做平，算你还清了国企的欠款。”乙答应并照办。对单位

% Source PDF physical page 399; see audit/source_coverage.json.
的应收账款（债权，财产性利益），甲有支配权，甲将其变成自已所有，让单位损失了应收账款，构成贪污罪，贪污数额是100万元，甲分得50万元，乙分得50万元好处。

（3）侵占的方式

由于侵占属于将自己已经占有的财物变成自己所有，而自己能够占有财物，是依照职权、职务占有财物。将自己依照职务占有的财物变成自己所有，表明渎职了，亵渎了职务。因此，侵占这种方式不要求再额外地利用职务便利。

例如，国有公司的出纳甲，并未使用其所保管的保险柜钥匙与密码，而是利用斧头劈开保险柜后取走现金，构成贪污罪，而不认定为盗窃罪（2008年第18题）。\footnote{这表明，“监守自盗”这个成语容易引起歧义。“监守”表明自己依照职务占有，将自己占有的财物变成自己所有，构成侵占，而非盗窃。所以，“自盗”这个表述不准确。此外注意，单位基于不法原因而让职员保管赃物。例如，单位收到一笔贿赂款，让出纳保管。出纳知道真相并保管。后出纳将该款据为己有，拒不归还单位。出纳构成掩饰、隐瞒犯罪所得罪。多数说认为，出纳不构成侵占罪，因为单位没有合法的返还请求权。基于此，出纳不构成侵占型贪污罪。但对这笔赃款应没收。不过，如果同事乙知道真相，从出纳处偷走这笔钱款，则乙构成盗窃罪。也即，对于乙而言，刑法保护单位对赃款的非法所有、保护出纳对赃款的非法占有。}

2.窃取

这里的“窃取”和盗窃罪的盗窃含义相同，也即，将他人占有的财物通过平和手段转移为自己占有，在这里就是指，将自己没有单独占有的公共财物，利用职务便利转移为自己占有。

[注意］贪污罪的盗窃手段，要求比普通盗窃罪多一个要件：利用职务便利。这里的利用职务便利，是指在对公共财物转移占有时，行为人的职权、职务发挥了实质贡献。

例1（2019年试题），某国企公司里，财务室的保险柜须有钥匙和密码共同使用才能打开，会计甲掌管钥匙，出纳乙掌管密码，甲乙均是国家工作人员。甲偷看乙的密码，打开保险柜，取走现金。保险柜的财物由甲乙共同保管、共同占有。共同占有的一方，将共同占有的财物转移为自己占有，属于盗窃行为。甲属于利用职务便利窃取公共财物，构成贪污罪。

例2，甲是国企的财务人员，乙是国企仓库管理员。甲盗窃仓库里的财物，构成盗窃罪。但是，甲无法在盗窃罪的基础上构成贪污罪。这是因为，甲完全没有占有财物，在盗窃时没有利用职务便利。

3.骗取

这里的“骗取”和诈骗罪的诈骗含义相同，也即，实施欺骗行为，使有处分权的人产生认识错误并处分财物，行为人因此取得财物。这就要求诈骗的公共财物在事实上不属于自己占有。

【注意】贪污罪的诈骗手段，要求比普通诈骗罪多一个要件：利用职务便利。这里的利用职务便利，是指在实施欺骗时，行为人的职权、权力发挥了实质贡献或影响力。受骗人之所以受骗并处分财物，是因为相信了行为人的职权、权力。

例1，国有建筑公司项目经理甲有权管理工人工资事宜，甲伪造工资表，会计乙基于对甲的职权行为的相信，只是形式审查，多发放了工人工资。甲构成贪污罪。

例2，某个科员甲出差，拿着假发票回到单位报销。单位的报销制度是，财务处对报销的发票和事项要实质审查其真实性，不是形式审查，不会轻易相信报账者。甲对此无权

% Source PDF physical page 400; see audit/source_coverage.json.
干涉，只能期待财务人员被骗。有次，财务人员被骗，给甲多报销了支出。由于甲在报销过程中不存在发挥职权施加影响力的问题，因此不构成贪污罪，只构成诈骗罪。

例3，投保人甲与国有保险公司理赔员乙（国家工作人员）勾结，乙利用工作身份欺骗上级领导丙，使丙处分了保险金。乙构成贪污罪，甲构成共犯。甲构成保险诈骗罪，乙构成共犯。

[提示1]实行行为不能缺。有些犯罪，具备了贪污罪的B1、B2、B3要件，但缺少A要件这个基础要件，也即缺少贪污的实行行为，不构成贪污罪。这一点请务必注意。例如，甲夜晚盗窃国有企业的仓库，先找到该企业的保安经理乙（国家工作人员），送一条烟，请其行个方便。乙便睁一只眼闭一只眼。甲盗窃既遂。乙不构成贪污罪，因为没有实施贪污罪的实行行为，只构成盗窃罪的帮助犯。

[提示2］变相贪污。模型：国企官员甲与私企老板乙勾结，伪造一个国企与私企的项目合同，甲让国企向私企打款，然后私企将款项打给甲。甲构成贪污罪的实行犯，乙构成贪污罪的帮助犯。

\subsubsection*{（四）主观目的}

贪污罪本质上是财产犯罪，要求具有非法占有公家财物的目的。该非法占有目的既包括为自己占有，也包括为第三人占有。

例1，官员甲利用职务便利，将公家财物送人，构成贪污罪。

例2，甲私营公司承揽乙国有公司的一个工程项目，乙国有公司的领导王某编造一个名目，基于此将乙公司的200万元打给甲公司。王某构成贪污罪。（2022年试题）

\subsection*{二、认定问题}

1.既遂。\CJKunderwave{判断标准}：就窃取、骗取而言，要求取得控制公共财物（建立自己的占有）；就侵吞而言，要求行使所有权。至于既遂后将财物又捐赠给公益事业，不影响既遂的成立。

[注意］犯罪既遂数额=公家损失数额。

例如（2014年主观题），国有化工厂副厂长甲（国家工作人员）与私营五金厂老板乙合谋，在化工厂的某贵重零件能够继续使用时，由甲利用职权，制造该零件报废、需要向五金厂购买的假象（该零件价值26万元），乙接到订单后，只向化工厂寄出供货单、发票而不需要实际供货，等化工厂将货款打到五金厂后，五金厂再将货款打到甲的账户。乙为了巴结甲，先将自己的26万元打到甲的账户，心想“反正化工厂的货款很快会打给我”。未料，甲的操作被人发现，化工厂未将“货款”打给五金厂。（1）甲构成贪污罪，虽然收到26万元，但不构成贪污罪既遂，因为化工厂没有损失，该钱是乙的钱。（2）乙不构成行贿罪，因为没有行贿意图，也即不是想将该钱赠送给甲，只是想临时垫付。乙构成甲贪污罪的帮助犯。

2.共犯与身份。

\InsertDiagram{figures/manual_table_p400_1.tex}

% Source PDF physical page 401; see audit/source_coverage.json.
\ContinuedTable{figures/manual_table_p401_1.tex}

例1，乙（民政局局长）为了贪污公款，让妻子甲帮助提供假发票。乙构成贪污罪（实行犯），甲构成贪污罪（帮助犯）。

例2-1，国有保险公司理赔经理乙（国家工作人员）和投保人甲内外勾结，各自分工。甲编造虚假索赔材料，提交给乙。乙提交给总经理丙。丙受骗，批出理赔款给甲。甲乙一起分掉该款项。（1）甲构成保险诈骗罪的实行犯。诈骗的实行行为的特点是，向受骗人传递虚假信息，意图使受骗人做出处分财物的行为。这个传递行为既可以是亲自传递，也可以是通过他人传递。甲通过乙向丙传递了虚假信息。乙构成甲的保险诈骗罪的帮助犯。（2）乙构成贪污罪的实行犯，属于欺骗型贪污，乙欺骗丙，实施了诈骗的实行行为。甲构成乙的贪污罪的帮助犯。（3）每个人均同时触犯两个罪，想象竞合，择一重罪论处。

例2-2，镇长乙负责土地征用补偿款的评估工作，与村民甲勾结。甲编造虚假申请材料，提交给乙，乙对此进行“评估”，将甲的虚假申请材料和自己的“评估”报告提交给负责审批的县长丙。丙受骗，批出补偿款给甲。甲乙一起分掉补偿款。（1）甲构成诈骗罪的实行犯。乙构成甲的诈骗罪的帮助犯（或共同正犯）。（2）乙构成贪污罪的实行犯。甲构成乙的贪污罪的帮助犯。（3）每个人均同时触犯两个罪，想象竞合，择一重罪论处。（2012年主观题）

例3-1，甲欲盗窃国有公司仓库货物，送给该公司保安经理乙（国家工作人员）一条烟，让“行个方便”。乙便不予阻止。甲盗窃既遂。乙构成甲的盗窃罪的帮助犯。乙不构成贪污罪，因为没有实施贪污罪的实行行为。

例3一2，甲为了骗领贫困补助，需要先成为“低保户”，为此找到自已的亲戚、民政局职员乙，乙负责审核“低保户”身份。乙明知甲不符合条件，仍授予甲“低保户”身份。甲凭借该身份，虚构申领材料，向社保局申领贫困补助3万元。社保局职员丙不知情，向甲发放了贫困补助。（1）甲实施了诈骗的实行行为，构成诈骗罪的实行犯，乙是诈骗罪的帮助犯（或共同正犯）。（2）乙不构成贪污罪的实行犯，因为乙的行为不属于贪污罪的实行行为。欺骗型贪污罪的实行行为，要求向受骗人传递虚假信息，意图使受骗人做出处分行为。乙只是授予甲一个假身份，乙自己不是申请人，也不是负责审核材料的人，乙自已没有意图使丙做出处分行为，乙只是帮助甲，甲意图使丙做出处分行为。注意，这一点与上述例2-1、例2-2不同。后两个案例中的乙是材料审核人和向处分人递交材料的人，对处分人丙实施了诈骗的实行行为。\footnote{可能有人认为，民政局职员乙的帮助行为起到关键作用，能够成为贪污罪的共同正犯。但是，共同正犯不能单独存在，必须依附正犯（实行犯）而存在。如果说乙构成贪污罪的共同正犯，其前提是存在贪污罪的正犯。但是，没有人构成贪污罪的正犯。}

% Source PDF physical page 402; see audit/source_coverage.json.
[巩固练习］例如，村民甲欲非法获取县民政局的危房补助款，与民政局局长乙串通，由甲负责虚构危房证明材料，上报民政局，乙明知真相，却将危房补助款批给甲。乙事后没有与甲分赃。乙实施了实行行为，构成贪污罪的实行犯，贪污罪的非法占有目的既包括为自己占有，也包括为第三人占有。甲是贪污罪的帮助犯。甲不构成诈骗罪，因为乙没有受骗。

3.罪名界限。

（1）贪污罪与职务侵占罪的区分，关键在于行为主体的身份：贪污罪的主体是国家工作人员；职务侵占罪的主体是非国家工作人员。

（2）贪污罪与私分国有资产罪的区分。私分国有资产罪（第396条第1款），是指国家机关、国有公司、企业、事业单位、人民团体违反国家规定，以单位名义将国有资产集体私分给个人，数额较大的行为。\CJKunderwave{两罪区分}：贪污罪是个人犯罪，私分国有资产罪是单位犯罪，为单位谋取非法利益。

4.处罚。根据第383条，犯贪污罪、受贿罪，被判处死刑缓期执行的，人民法院可以同时决定在其死刑缓期执行二年期满依法减为无期徒刑后，\CJKunderwave{终身监禁，不得减刑、假释}。

\section{挪用公款罪}

[问题】官员甲利用职务便利，挪出公款100万元，以本单位的名义，借给乙单位使用，私下收了乙单位5万元。对甲如何处理？（挪用公款案）

第384条国家工作人员\CJKunderwave{利用职务上的便利}，挪用公款\CJKunderwave{归个人使用}，\CJKunderwave{进行非法活动}的，或者挪用公款\CJKunderwave{数额较大、进行营利活动}的，或者挪用公款\CJKunderwave{数额较大、超过三个月未还}的，是挪用公款罪，处五年以下有期徒刑或者拘役；情节严重的，处五年以上有期徒刑。\CJKunderwave{挪用公款数额巨大不退还的}，处十年以上有期徒刑或者无期徒刑。

挪用用于救灾、抢险、防汛、优抚、扶贫、移民、救济款物\CJKunderwave{归个人使用}的，\CJKunderwave{从重}\CJKunderwave{处罚}。

挪用公款罪，行为主体是国家工作人员。\footnote{提示：根据第382条第2款规定，受国家机关、国有公司、企业、事业单位、人民团体委托管理、经营国有财产的人员，可以成为贪污罪的主体。这类人本不属于国家工作人员，但是该款将他们拟制为国家工作人员。由于第384条挪用公款罪中没有类似法律拟制，因此此类人挪用受委托管理、经营的国有财产的，只能定挪用资金罪，不能定挪用公款罪。}行为对象是公款，也即公家的资金。

\subsection*{一、行为公式}

\InsertDiagram{figures/manual_p402_1.tex}

本罪的实行行为（构成要件行为）只有挪出行为，不包括使用行为。这是因为，如果认为本罪的实行行为由挪出行为和使用行为组成，那么会带来问题。

例如，甲挪出公款进行贩卖毒品，应定挪用公款罪和贩卖毒品罪，数罪并罚；但是由于使用行为（贩卖毒品）是挪用公款罪的实行行为，对两罪并罚，就会导致使用行为被评

% Source PDF physical page 403; see audit/source_coverage.json.
价了（处罚了）两次，违反了禁止重复评价原则。但是，又必须对两罪并罚，为此，只好将使用行为逐出挪用公款罪的实行行为范围，使用行为便成为实际处罚条件。亦即，虽然使用行为不是本罪的实行行为，但只有具备了使用行为，才能对行为人实际启动刑罚。\footnote{上述图表中，“归个人使用行为”在外围打个虚框，意思是使用行为不是本罪的实行行为。在内圈打个实框，意思是使用行为是处罚条件。}

[结论】本罪的实行行为是，为了归个人使用，挪出公款的行为。只要实施了该实行行为，就成立本罪。如果因为意志以外原因未能挪出来，就成立犯罪未遂。如果挪出来，就成立犯罪既遂。本罪的特殊性在于，即使本罪既遂了，并不直接科处刑罚，还要求具备实际处罚条件（归个人使用的行为）才予以处罚，追究刑事责任。如果不具备实际处罚条件，就不追究刑事责任，从这个意义上讲，也就不构成犯罪，也即实际处罚条件也属于广义的犯罪成立条件。\footnote{有的同学可能困惑：“挪出来就犯罪既遂了，为何又说不构成犯罪？”这是因为，“不构成犯罪”需要在不同阶段来理解。例如，某女明星逃税数额特别巨大，已经构成逃税罪既遂，但是由于补缴税款，不予追究刑事责任。从最终是否追究刑事责任来讲，她不构成犯罪，但她前面的确构成逃税罪既遂。}因此，如果挪出来，但未进行违法活动，或三个月内还了，不能定犯罪中止，而是不处罚、不定罪。（2011年第54题）

[巩固练习]题1，甲为了贩毒而挪出公款后，乙参与贩毒。乙不构成挪用公款罪的承继的共同犯罪，乙构成贩卖毒品罪的共同犯罪。甲构成挪用公款罪与贩卖毒品罪，并罚。

题2，乙为了贩毒，指使甲挪出公款，甲明知乙使用公款会用于贩毒，而将公款挪用给乙。乙拿到公款后进行贩毒。甲成立挪用公款罪的实行犯，对乙的贩毒行为也要负责，构成贩卖毒品罪的帮助犯，数罪并罚。乙构成挪用公款罪的教唆犯和贩卖毒品罪的实行犯，并罚。

\subsection*{二、“归个人使用”}

\subsubsection*{（一）行为定性}

挪用公款罪，要求将公款归个人使用，也即挪作私用。如果挪作公用，不构成挪用公款罪。根据立法解释、司法解释，归个人使用的情形有三种：\footnote{参见2002年4月28日全国人民代表大会常务委员会《关于\ensuremath{<}中华人民共和国刑法\ensuremath{>}第三百八十四条第一款的解释》；2003年11月13日《全国法院审理经济犯罪案件工作座谈会纪要》；2026年5月1日最高人民法院、最高人民检察院《关于办理贪污贿赂刑事案件适用法律若干问题的解释（二）》。}

第一，将公款归本人、亲友或其他自然人使用。（个人落了人情）

第二，以个人名义将公款供其他单位使用。（个人落了人情）

第三，个人决定以单位名义将公款供其他单位使用，谋取个人利益。（个人落了好处）\footnote{第一，其中的“个人利益”，既包括不正当利益，也包括正当利益；既包括财产性利益，也包括非财产性利益，但这种非财产性利益应当是具体的实际利益，如升学、就业等。第二，“个人利益”，并不限于一个人的利益，而是相对于单位、集体而言，为单位少数特定人谋取利益，也属于谋取个人利益。第三，“谋取个人利益”，既包括行为人与使用人事先约定谋取个人利益的情况，也包括虽未事先约定但实际已获取了个人利益的情况。}（2010年第20题）

例如，本节开头“挪用公款案”中，甲属于挪作私用，构成挪用公款罪和受贿罪，并罚。

【注意】上述第二种情形包括：个人通过虚构付款事由或者将单位应收账款不按规定人账等逃避单位监管的方式，将公款提供给其他单位使用。

例1，甲系某国企财务总监，朋友乙的丙公司资金周转困难。甲未经单位集体研究、

% Source PDF physical page 404; see audit/source_coverage.json.
未履行借款审批流程，虚构“该国企向丙公司采购原材料”的付款事由，伪造采购合同，从国企对公账户转账500万元至丙公司，供丙公司用于经营周转，无真实采购业务。甲虚构付款事由、逃避单位监管，虽为公对公转账，仍属于“以个人名义将公款供其他单位使用”，并且用于营利活动，构成挪用公款罪。

例2，甲系某国企负责人，该国企对乙私营公司有应收账款300万元。乙公司要还款，甲逃避单位监管，让乙公司将款项直接转给丙公司（甲的朋友公司），用于营利活动。甲构成本罪。

这表明，挪用公款罪的公款包括：\CJKunderwave{一是单位公账上的现有公款；二是单位的应收}\CJKunderwave{账款}。

\subsubsection*{（二）行为方式}

1.内容

（1）“进行非法活动”。

第一，不要求挪用时间达3个月。法条对挪用数额未作要求，司法解释要求是3万元。

第二，非法活动如果构成其他犯罪，应数罪并罚。

（2）“数额较大，进行营利活动”。不要求挪用时间达3个月。司法解释规定，“数额较大”是5万元。

（3）“数额较大，超过3个月未还”。这是指进行非法活动、营利活动以外的一般活动，在时间上超过3个月未还。\footnote{如果不进行一般活动，而是放家里不使用，只要超过3个月未还，也构成本罪。此时，客观处罚条件不是使用行为，而是使公款脱离单位控制长达3个月之久。}司法解释规定，“数额较大”是5万元。

2.包容评价

三种实际处罚条件的包容评价。根据法益危害程度，非法活动可以包容评价为营利活动，营利活动可以包容评价为一般活动。例如，甲挪用公款2万元进行营利活动，挪用公款3万元进行其他活动，均超过3个月未还。如果单独看，两个行为都不构成本罪，但是，将2万元的营利活动评价为一般活动，则一般活动为5万元，符合“数额较大，超过3个月未还”，成立本罪。\footnote{练习：题1，甲挪用公款2万元进行赌博、走私等非法活动，挪用公款3万元进行营利活动。如果单独看，两个行为都不构成本罪，但是，将2万元的非法活动评价为营利活动，则符合“数额较大（5万元），进行营利活动”，成立本罪。题2，甲挪用公款1万元进行非法活动，3个月内归还；挪用公款2万元进行营利活动，3个月内归还；挪用公款2万元进行其他活动，超过3个月未还。申不构成本罪，因为非法活动、营利活动包容评价为一般活动，要求超过3个月未还。}（2014年第62题）

3.认定

（1）在三种实际处罚条件内部，可以转换。

例1，甲为了走私而挪出公款，但因为股市行情好就用于炒股。对此应认定为挪用公款进行营利活动。甲成立本罪。

例2，甲为了炒股而挪出公款（数额较大），但因为股市行情不好就一直将公款放着没用，超过3个月未还。对此就认定为挪用公款数额较大，超过3个月未还。甲成立本罪。

（2）三种实际处罚条件是故意私用行为，所以也需要遵守主客观相一致原则。主观上

% Source PDF physical page 405; see audit/source_coverage.json.
没有想实施三种实际处罚条件，客观上实施了，最终不构成本罪。\footnote{所以，该实际处罚条件与一般的客观处罚条件有所不同。后者只看客观上是否具备。另外需注意，实际处罚条件、客观处罚条件，都不是四个要件内部的要素，而是四要件体系之外、之后的条件。}

例如，乙谎称为了看病需要钱，唆使甲挪用公款30万元给自己，两周后归还。甲信以为真并照办。乙实际上拿着公款去贩卖毒品，两周后归还给了甲。（1）甲挪出公款，就构成既遂，接下来看是否满足实际处罚条件。甲客观上实施了处罚条件“进行非法活动”，但主观上没有实施该处罚条件的意图。所以，甲最终不构成挪用公款罪。不过，由于甲客观上实施了处罚条件“进行非法活动”，所以，在处罚条件的客观层面，甲构成“挪用公款罪”，是个暂时的“犯罪”概念。（2）乙虽然欺骗了甲，但不构成挪用公款罪的间接正犯，因为成立挪用公款罪的正犯要求具有国家工作人员的身份，乙没有国家工作人员的身份。不过，乙可以构成挪用公款罪的教唆犯。甲属于处罚条件的客观层面的“实行犯”。乙构成挪用公款罪的教唆犯和贩卖毒品罪，并罚。甲不构成贩卖毒品罪的帮助犯，因为主观上没有帮助故意。\footnote{有人认为，由于“进行非法活动”可以评价为“进行一般活动”，所以，甲不构成处罚条件的客观层面的“实行犯”。然而，包容评价的做法主要是为了给行为人定罪，而非出罪或判处轻罪。例如，甲故意杀人，不能因为故意杀人罪可以包容评价为故意伤害罪，就给甲定故意伤害罪。}

\subsection*{三、挪用公款罪与贪污罪的关系}

挪用公款罪是A，贪污罪是A+B，B就是非法占有目的。二者的区分标准是有无非法占有目的（有无归还意思）。两罪是包容评价关系，贪污罪可以包容评价为挪用公款罪。当无法查明非法占有目的时，可认定为挪用公款罪。（2016年第11题）

根据司法解释，有下列情形就表明行为人具有非法占有的目的：

1.携带挪用的公款潜逃的。对其携带潜逃的公款部分，定贪污罪。注意罪数问题。

例1，甲挪用100万元公款，已构成挪用公款罪。甲又携带该100万元潜逃，构成贪污罪。贪污罪吸收挪用公款罪，定贪污罪。

例2，甲挪用100万元公款，已构成挪用公款罪。甲又携带其中的80万元潜逃，构成贪污罪，数额是80万元。此时，挪用公款罪的数额是100万元-80万元=20万元。对挪用公款罪（20万元）和贪污罪（80万元）并罚。

例3，甲挪用单位A账户的公款，构成挪用公款罪。甲又贪污单位B账户的公款，构成贪污罪。由于行为对象不同，对甲的挪用公款罪和贪污罪应并罚。

2.采取虚假发票平账、销毁账目等手段，使所挪用的公款在账上难以反映，且不归还的。（2008年第92题）

3.截取单位收入不人账，非法占有，使所占有的公款在账上难以反映，且不归还的。

4.有能力归还公款而拒不归还。这表明，贪污罪是因主观原因（不想还），给公家造成损失。挪用公款罪是因为客观原因，给公家造成损失，也即主观想还，客观还不了。例如，官员王某挪用公款100万元炒期货，完全亏损。王某深感无力还债，便潜逃外地。王某构成挪用公款罪。（2022年试题）

% Source PDF physical page 406; see audit/source_coverage.json.
\subsection*{四、挪用型犯罪的总结}

\InsertDiagram{figures/manual_table_p406_1.tex}
\SourceFootnote{\textbf{表注①}\quad 少数说认为，挪用特定款物罪是单位犯罪。多数说认为，法条没有规定该罪是单位犯罪，并且没有处罚单位，所以该罪是自然人犯罪。不过，该罪的行为人不能将特定款物挪作个人使用，否则就定挪用公款罪。}

例如，官员王某将救灾款转用于装修单位办公大楼，构成挪用特定款物罪。（2022年试题)

\section{受贿罪}

[问题］老板乙有求于官员甲，某个暑假，安排官员甲全家去欧洲旅游两星期，所有花费（20万元）均由乙支出。甲是否构成受贿罪？（免费旅游案）

第385条国家工作人员\CJKunderwave{利用职务上的便利}，\CJKunderwave{索取他人财物的}，或者\CJKunderwave{非法收受他人财}\CJKunderwave{物}，\CJKunderwave{为他人谋取利益的}，是受贿罪。

国家工作人员在经济往来中，违反国家规定，收受各种名义的回扣、手续费，归个人所有的，以受贿论处。

\subsection*{一、保护法益}

例如（打信息差案），某高考考生的成绩完全可以选择某重点高校的任何专业。但是，考生家长不知道这一情况，希望招生人员甲能够满足孩子选择某个专业的愿望。甲趁机说：“只要给我5万元，我就能帮忙办成此事。”家长便给了甲5万元。实际上，甲没有做任何违规的事情。该考生选到了自已想要的专业。

少数说认为，受贿罪的保护法益是职务行为的公正性。由于甲没有做任何违规的事情，没有侵害职务行为的公正性，因此不构成受贿罪。甲仅构成诈骗罪。

多数说认为，\CJKunderwave{受贿罪的保护法益是职务行为的不可收买性}。甲许诺“只要给钱，就办事”，会让家长认识到，甲的职务行为是可以被收买的。甲便侵害了“职务行为的不可收

% Source PDF physical page 407; see audit/source_coverage.json.
买性”这个法益。因此，甲构成受贿罪。甲同时构成诈骗罪，与受贿罪想象竞合，择一重罪论处。\footnote{家长不构成行贿罪，因为家长谋取的利益在客观上具有正当性。对此，参见后文行贿罪的讲解。}

\subsection*{二、实行行为之交易标的}

受贿罪的实行行为是交易行为（权钱交易），用钱买权，用权收钱，二者形成交易关系（对价关系、条件关系）。卖方是国家工作人员，出卖的是服务（职务行为）；买方是请托人，用钱财购买职务行为。

第一，国家工作人员单纯的收钱行为，不是受贿罪的实行行为。例如，官员儿子结婚，老板赠送巨额礼金。如果查不出来，该礼金是官员某个职务行为的对价，则不构成受贿罪。

第二，国家工作人员单纯的办事行为，不是受贿罪的实行行为。例如，官员给老板办事时没有收钱意图（受贿故意），办事后索取报酬，老板给钱。基于已办成的事来收钱，形成对价关系，成立受贿罪。有人认为，办事是受贿罪的实行行为，由于办事时没有受贿故意，行为与故意不是同时存在，因此不构成受贿罪。实际上，办事行为不是受贿罪的实行行为。

\InsertDiagram{figures/manual_p407_1.tex}

请托人的财物（贿赂）

\subsubsection*{（一）}

贿赂，是指职务行为的不正当报酬。人罪数额是3万元。贿赂只包括财物和财产性利益，但不包括劳务本身，也不包括帮国家工作人员办件不存在财产性利益的事情。

司法解释\footnote{2016年4月18日最高人民法院、最高人民检察院《关于办理贪污贿赂刑事案件适用法律若干问题的解释》。}规定：贿赂犯罪中的“财物”，包括货币、物品和财产性利益。财产性利益包括可以折算为货币的物质利益如房屋装修、债务免除等，还包括需要支付货币的其他利益如会员服务、旅游等。后者的犯罪数额，以实际支付或者应当支付的数额计算。

例1，为官员提供房屋装修、出国旅游等，存在财产性利益。本节开头“免费旅游案”中，甲构成受贿罪。

例2，官员在色情场所，由请托人支付费用的，或者请托人支付费用雇请卖淫者为官员提供性服务的，存在财产性利益，属于受贿。但是，请托人自己直接为官员提供性服务的，不存在财产性利益，不是受贿。

例3，为官员写论文、作业，不是受贿。

例4，为官员谋取晋升机会，为官员的子女找工作，不是受贿。

\subsubsection*{（二）自己的职务行为（办事）}

国家工作人员利用自己的职务行为（利用职务便利）换取请托人的财物。

1.“职务”。这里的职务是指国家工作人员的公务。

例1，甲是公立大学科研处处长（国家工作人员），在业余时间，利用自已的知识为某公司研发产品，获取报酬，不构成受贿罪，无罪。

% Source PDF physical page 408; see audit/source_coverage.json.
例2，甲是公立医院副院长（国家工作人员），收受医药代表5万元，承诺在坐诊看病开处方时多开该医药公司的药品。坐诊看病是纯粹技术劳动，不具有行政职责性，不是公务活动，因此不构成受贿罪，但构成非国家工作人员受贿罪，也即此时甲的身份性质是非国家工作人员。如果甲承诺“本医院采购科会多采购该医药公司的药品”，则甲构成受贿罪。

2.“职务便利”。这包括两种类型。一是直接型。这是指利用职务本身直接为请托人谋取利益。例如，主管市政工程的副市长甲收老板乙的钱，将某个市政工程交给乙。二是间接型。这是指，利用职务的影响力，让第三人为请托人谋取利益。这里要求国家工作人员对第三人的影响力达到制约性的程度。

例1，老板丙的民营企业依靠副市长甲发展壮大。老板乙想承包丙的企业的一个项目，给甲送钱，甲指示丙“配合一下”，丙照办。甲构成受贿罪。

例2，老板乙想找某市税务局局长丙“办事”，给该市副市长甲送钱，甲指示丙“配合一下”，丙照办。甲构成受贿罪，不属于后文的“斡旋受贿”。

3.“利用”。行为人利用自己的职务行为作为筹码，与他人的财物形成对价关系（交易关系）。这是一种“无A则无B”的必要条件关系，也即，请托人如果不是看中行为人的职务行为，是不会给钱的。如果“办事”与“收钱”之间，由于时间长、数量小等原因，导致不能或者难以认定存在对价关系，则不认定为受贿。

4.利用职务行为，为请托人谋取利益（办事），既包括正当利益，也包括不正当利益。

5.利用职务行为，为请托人谋取利益（办事），所谋取的利益的受益人，既可以是请托人本人，也可以是请托人指示的第三人，均属于为请托人办事。例如，乙欲获得晋升，又不宜直接向大领导丙行贿，于是将实情告知商人朋友甲，甲提出由其向丙行贿并请托。丙接受财物，并为乙办了事。丙构成受贿罪。甲乙成立行贿罪的共犯。

\subsection*{三、实行行为之交易方式}

\subsubsection*{（一）事前型交易（事前受贿)}

这是指办事之前就有交易约定。

1.索取贿赂

受贿罪的成立条件：只要实施了索贿行为，就成立受贿罪。在此，不要求官员许诺办事（为他人谋取利益）。

受贿罪的既遂条件：要到了钱。至于官员有无实际办事（为他人谋取利益），不重要，不影响受贿罪的成立与既遂。

2.收受贿赂

（1）受贿罪的成立条件：请托人提出送钱，官员许诺办事。

许诺是一种意思表示，既可以是明示的，也可以是暗示的。司法解释规定：①明知他人有具体请托事项，给予财物时不拒绝，就是暗示许诺为他人谋取利益。②国家工作人员索取、收受具有上下级关系的下属或者具有行政管理关系的被管理人员的财物价值3万元以上，可能影响职权行使的，视为承诺为他人谋取利益。\footnote{2016年4月18日最高人民法院、最高人民检察院《关于办理贪污贿赂刑事案件适用法律若干问题的解释》。}（2017年第62题）

% Source PDF physical page 409; see audit/source_coverage.json.
（2）受贿罪的既遂条件：官员收了钱。

第一，至于官员有无实际办事（为他人谋取利益），不重要，不影响受贿罪的成立与既遂。（2007年第65题）

第二，官员收受请托人的钱财，既可以亲自收，也可以让请托人交给第三人。让请托人交给第三人，既包括让第三人转交给自己（官员），也包括让第三人自己占有。也即，官员对贿赂的非法占有目的，既包括为自己非法占有，也包括为第三人非法占有。

[结论］受贿罪的收钱，既包括为自己收钱，也包括为第三人收钱。

第三，第三人明知请托人提供的是贿赂，成立受贿罪的共犯；不明知，则不成立受贿罪共犯。例如，老板甲要给官员乙送钱，乙让甲送给丙（乙的情妇），甲照办。乙让丙自已留用。乙构成受贿罪。若丙知道钱的来路，则成立受贿罪的共犯。

\subsubsection*{（二）事后型交易（事后受贿）}

这是指，国家工作人员在办事时，没有收受财物的意图，办事之后产生了收受财物的意图，将自己已经办的事作为筹码进行交易，获取不正当报酬。

【结论】这种行为构成受贿罪。司法解释规定：履职时未被请托，但事后基于该履职事由收受他人财物的，属于“为他人谋取利益”，构成受贿罪。\footnote{2016年4月18日最高人民法院、最高人民检察院《关于办理贪污贿赂刑事案件适用法律若干问题的解释》。}

例如，税务局局长甲没有受贿意图，为乙的公司办理了正当业务。事后，乙欲感谢甲，甲说：“我办事时没想图你们的钱。”乙说：“正因如此，更要感谢您！”便送甲4万元，甲予以接受。甲成立受贿罪。乙不构成行贿罪，因为乙谋取的是正当利益。\footnote{早前观点认为，事后型交易不成立受贿罪，理由是，成立受贿罪要求遵守“行为与故意同时存在原则”，而国家工作人员办事时没有受贿故意。这种看法是错误的。受贿罪的实行行为不是实施职务行为为他人办事，而是凭借职务行为与他人的财物进行交易，交易行为才是受贿罪的实行行为。凭借已经办的事与他人的财物交易时，国家工作人员具有受贿故意，因此成立受贿罪，在此并不违反“行为与故意同时存在原则”。这一点，得到了现在司法解释的认可。}

[注意】离职问题

（1）在职时，约定“先办事，离职后再收钱”，构成受贿罪。这是因为，约定就是达成交易，达成交易就是受贿罪的实行行为。实施该行为时，行为人具有国家工作人员身份，因此成立受贿罪。不能认为只有收钱行为才是受贿罪的实行行为。如果离职后没收到钱，则构成受贿罪未遂。

（2）在职时，没有约定，离职后收钱，不构成受贿罪。这是因为，此时的收钱是受贿罪的实行行为，但是此时行为人已经不具有国家工作人员身份了。（2003年第38题，2007年第65题，2009年第64题）

\subsubsection*{（三）司法解释规定的变相受贿\protect\footnotemark}\footnotetext{2007年7月8日最高人民法院、最高人民检察院《关于办理受贿刑事案件适用法律若干问题的意见》；2026年5月1日最高人民法院、最高人民检察院《关于办理贪污贿赂刑事案件适用法律若干问题的解释（二）》。}

1.以交易形式收受贿赂的，成立受贿罪。例如，房地产公司老板甲求官员乙办事，向乙出售房屋。市场价是200万元，甲以50万元卖给乙。甲构成行贿罪，乙构成受贿罪。

2.收受干股（未实际出资而获得的股份）的，成立受贿罪。

% Source PDF physical page 410; see audit/source_coverage.json.
3.以开办公司等“合作”投资名义，未实际出资，以此收受贿赂的，成立受贿罪。受贿数额为出资额。

4.以委托请托人投资证券、期货或者其他委托理财的名义，未实际出资而获取“收益”，或者虽然实际出资，但获取“收益”明显高于出资应得收益，以此收受贿赂的，成立受贿罪。受贿数额为获取“收益”的数额。

5.以赌博形式输给官员，以此收受贿赂的，成立受贿罪。

6.以特定关系人“挂名”领取薪酬的，成立受贿罪。

7.名为借用，实为受贿，成立受贿罪。真借用还是假借用，判断因素：有无借用的合理事由、是否实际使用、借用时间长短、有无归还条件等。

8.行贿人按照受贿人授意购买特定物品后给予受贿人的，应当以行贿人实际支付的购买金额认定受贿数额。

例1，甲找官员乙办违规事项，乙让甲到丙字画店买一幅画，该幅画的市场价是2千元，甲以5万元购买，送给乙。乙再将该幅画交给丙字画店，丙字画店将5万元交给乙。（清代官场的一种受贿方式）。甲构成行贿罪，数额是5万元。乙构成受贿罪，数额是5万元。

例2，甲找官员乙办违规事项。乙是二流画家，画了一幅画，市场价是2千元。甲以5万元从乙手里购买。甲构成行贿罪，实际行贿数额是4.8万元，这是因为虽然甲支付了5万元，但也收到了价值2千元的财物。乙构成受贿罪，实际受贿数额是4.8万元。这是因为，乙虽然收到5万元，但也支付了2千元的财物。

\subsection*{四、主观故意}

1.国家工作人员主观上没有受贿罪故意，客观上收受财物后及时退还、上交。这种行为不构成受贿罪。

例1，老板甲将银行卡塞进官员乙家的沙发缝隙里，乙1个月后发现，及时上交。

例2，老板甲送官员乙一盆兰花，声称价值只有100元。2个月后乙得知该兰花价值5万元，及时上交。

［提示］国家工作人员受贿后，因自身或者与其受贿有关联的人、事被查处，为掩饰犯罪而退还或者上交的，依然构成受贿罪。

2.数额上的主客观相一致

受贿罪的第一档法定刑条件是“数额较大（3万元）”，第二档法定刑条件是“数额巨大（20万元）”，第三档法定刑条件是“数额特别巨大（300万元）”。

基本原理：由于受贿罪是故意犯罪，作为收受的对象，\CJKunderwave{“数额较大（巨大、特别巨}\CJKunderwave{大）的财物”是个主客观相统一的要素}。\CJKunderwave{常见误解是，将“数额较大（巨大、特别巨大）}\CJKunderwave{的财物”理解成纯粹客观要素，只看客观上是否具备}。由此导致如下误解。

% Source PDF physical page 411; see audit/source_coverage.json.
\InsertDiagram{figures/manual_table_p411_1.tex}

\subsection*{五、既遂问题}

\subsubsection*{（一）既遂标准}

既遂标准：国家工作人员带着受贿故意，接受了他人财物，取得控制了财物，就是受贿罪的既遂。

1.受贿罪的既遂，只要求在事实上取得控制财物，不要求在民法上取得财物的所有权。

例如，甲送官员乙一套房，虽然没有过户，但只要乙收下，就构成既遂。

[巩固练习]题1（2022年试题），官员乙收到甲的银行卡，卡内有资金，设有密码，乙不知密码，无法支配卡内资金。乙构成受贿罪未遂。

题2（2022年试题），官员丙收到甲的银行卡，卡内有资金，丙知道密码，但卡内资金是定期存款，不是活期存款。丙构成受贿罪既遂，因为虽然是定期存款，但是丙已经事实上取得控制。

% Source PDF physical page 412; see audit/source_coverage.json.
2.既遂后，又被行贿人取走的，不影响既遂结论。

例如（2014年第54题），乙用自己的身份证办了一张银行卡，存入50万元，行贿给官员甲，并告知密码。甲收下该银行卡。后乙得知甲被免职，将该卡挂失取回50万元。甲收下卡就是受贿罪既遂。虽然从名义上看，乙仍是卡的所有权人，但甲构成既遂，也即甲在事实上占有了卡里资金。乙构成行贿罪既遂，表明乙在事实上对卡里资金失去了占有。多数说认为，乙的挂失行为导致甲对卡里资金失去占有，乙建立了自已的占有，构成盗窃罪。乙前面行为构成行贿罪既遂，后面行为构成盗窃罪既遂，数罪并罚。少数说认为，乙是卡的法律上的名义人，乙挂失自已的卡，取自己的钱，不构成盗窃罪。

[巩固练习]题1（2023年试题），乙用自己的身份证办了一张银行卡，存入100万元，行贿给官员甲，并告知密码。甲收下。次日，乙为帮在银行工作的朋友完成业绩，又将该卡中的100万元资金全部转为定期存款，等朋友的事情结束，乙会转回活期存款。在此期间，甲无法使用该笔资金。甲收下卡，就构成既遂，此后不能使用，不影响既遂。乙只是临时使用一下，具有归还意图，没有非法占有目的，因此不构成盗窃罪。

题2（2024年试题），乙向官员甲行贿，乙设立一个炒股账户，打入500万元，将账户密码交给甲。日后甲用这笔钱炒股赚了80万元。后甲未办成乙所托之事，乙便修改密码，甲无法取钱。乙将存有500万元的账户密码交给甲，甲构成受贿罪既遂，数额是500万元，80万元是犯罪所得的收益。乙构成行贿罪既遂。此后乙修改密码，导致甲对账户失去控制，由乙控制账户。乙构成盗窃罪。

3.既遂后，如何处置，不影响既遂结论。（2026年试题）

例1，官员收受了他人的转账支票，没有提取现金的；收受购物卡后，没有购物的；收受银行卡，没有使用的，都不影响既遂结论。（2006年第19题，2015年第5题，2022年试题）

例2，官员收受贿赂后，将贿赂用于公家事务、公益事业，不影响既遂结论。

例3，官员收受贿赂后，向对方回赠财物的，不能从受贿数额中扣除回赠的数额。

4.两人共同受贿，受贿的数额是财物整体的数额，而非两人各自分得的数额。例如，甲请国税局两位副局长乙丙共同解决漏税问题，乙丙收取10万元，乙分得4万元，丙分得6万元。乙丙的受贿数额都是10万元。

5.时间标准

既遂数额（受贿数额）是指犯罪所得数额，以犯罪时为准。犯罪所得的收益的数额，以案发时为准。例如，老板甲向局长乙行贿，送乙一套房，根据当时市场行情价，价值100万元，但当时没有办理产权过户。乙将房屋用于出租，出租了一年，获得房屋租金10万元。此后，双方办理了房屋过户手续，过户到乙名下。此时该房屋的市场价已经是120万元，此时案发。根据行为与行为对象同时存在原则，受贿数额是100万元。租金10万元和涨价的20万元属于犯罪所得的收益。（2026年试题）

\subsubsection*{（二）涉及股权的问题总结}

1.收受干股

这是指，不用实际出资，获得股份（股权）。在此，受贿罪的对象（贿赂）是指股权本身，不是股权的预期收益。

（1）成立受贿罪，不要求股权登记转让给受贿人，可以由行贿人代持。

% Source PDF physical page 413; see audit/source_coverage.json.
（2）未转让登记的，受贿数额以实际分红计算。

（3）转让登记的，受贿数额以转让时股份市场价值计算，分红属于受贿所得的收益。

（4）双方约定数额达到3万元，但以上受贿数额未能达到3万元的，定受贿罪未遂。（2020年试题、2024年试题）

2.股权的预期收益

这是指，官员出资购买了股权，所以股权不是贿赂，贿赂是指该股权的预期收益。该预期收益应具有可期待性，也即较大的实现可能性。

司法解释规定：以收受股票、股权的预期收益作为贿赂形式，构成犯罪的，受贿数额按照案发时实际获利认定；案发时尚未实际获利的，受贿数额一般按照案发时涉案资产的市场价格与支付价格的溢价认定。\footnote{2026年5月1日最高人民法院、最高人民检察院《关于办理贪污贿赂刑事案件适用法律若干问题的解释

（二）》。}

（1）针对股权的预期收益

例如（原始股），乙请托局长甲帮忙办违规事项，乙让甲获得购买乙公司原始股的资格，一般人无法获得该资格，乙公司即将上市，一旦上市，股价大概率会大涨；甲出资100万元，购买乙公司的原始股。乙公司上市后，甲的股价涨到600万元，甲卖出股票，实际获利500万元。甲的受贿数额是500万元。如果案发时，甲的股价仍是600万元，甲没有卖出股票，那么受贿数额仍是500万元。如果乙公司上市后，未料股价没涨或大跌，那么甲没有获利，受贿数额是零，构成受贿罪未遂。（2021年试题）

（2）针对股票的预期收益

这里的股票不是指普通二级市场股票，而是指上市公司非公开发行的股票，也即定向增发的股票。一般人没有资格购买。

例1，乙是上市公司实际控制人，请托局长甲帮忙办违规事项，专门给甲一个非公开发行名额，给甲的定增认购价大幅低于市场价。甲出资100万元，购买该定增股。案发时，甲以600万元卖出，实际获利500万元。甲构成受贿罪，受贿数额是500万元。

例2（股票本身作为贿赂），老板乙请局长甲办违规事项，给甲赠送市值100万元的股票，做了转让登记。在此，受贿罪的对象（贿赂）就是该股票本身，受贿数额（犯罪所得数额）是100万元。案发时，该股票溢价了50万元，该50万元属于犯罪所得收益。（2026年试题）

\subsubsection*{（三）违禁品的价值数额}

违禁品（假币、毒品、淫秽物品）具有财物价值（消极价值），在不法分子手里能够卖钱。在计算违禁品的价值数额时，可参考其非法交易的价值数额（行情价）。\footnote{有人可能认为，对违禁品可以按照情节轻重来论，不需要计算价值数额。但是，第一，如何判断情节轻重，需要考虑价值数额。例如，官员收受10公斤毒品和收受10克毒品，在情节上必然有区别。第二，受贿罪的既遂标准是收受到数额较大（3万元）的财物。为了判断是否既遂，需要衡量违禁品的价值大小。}

例1，老板答应送给官员2公斤海洛因，官员收到。经过评估，2公斤海洛因价值10万元（非法交易的行情价）。官员构成受贿罪既遂。

例2，老板答应送给官员50万元，将大袋子塞给官员，然后离去。官员打开发现是50万元假币。经过评估，50万元假币的价值有3万元（非法交易的行情价）。官员构成受

% Source PDF physical page 414; see audit/source_coverage.json.
贿罪既遂。如果官员打开发现只有1000元假币，则构成受贿罪未遂。（2009年第60题）

例3，老板答应送给官员两张淫秽光盘，官员收到。官员不构成受贿罪。

\subsection*{六、罪数问题}

\subsubsection*{（一）受贿罪与渎职罪的关系}

行为模型：收钱（构成受贿罪）+办事（构成渎职罪）

由于成立受贿罪，不要求官员为请托人实际办事（为他人谋取了利益），因此，实施办事行为不是受贿罪的构成要件行为（实行行为）。基于此，实施办事行为若构成其他犯罪（主要是渎职罪），与受贿罪属于两个行为，构成两个罪，原则上应数罪并罚。\CJKunderwave{总结}：受贿罪+办事行为（渎职罪）=数罪并罚。

例1（2013年第56题），官员甲收受境外组织的钱，将国家秘密提供给该组织，甲构成受贿罪和为境外非法提供国家秘密罪，数罪并罚。

例2，官员甲收受老板乙的钱，为乙的公司徇私舞不征税款，甲构成受贿罪与徇私舞弊不征税款罪，数罪并罚。

例3，狱警甲答应犯人家属请求，私自释放犯人，收取家属钱财，甲构成私放在押人员罪和受贿罪，数罪并罚。

例4，官员甲收受乙的钱，为乙挪用公款，用于非法活动。甲构成受贿罪和挪用公款罪，数罪并罚。

[例外］受贿罪+徇私枉法罪等四罪=牵连犯，择一重罪论处。

第399条第4款司法工作人员收受贿赂，有前三款行为的（徇私枉法罪，民事、行政枉法裁判罪，执行判决、裁定失职罪，执行判决、裁定滥用职权罪），同时又构成受贿罪的，依照处罚较重的规定定罪处罚。

其中的“收受贿赂”既包括收受贿赂，也包括索取贿赂。

\subsubsection*{（二）受贿罪与贪污罪的关系}

1．二者的区分

区分标准：看财物的来源属性。受贿罪收受的财物是行贿人的个人钱财。贪污罪贪污的是公共财物。

例1，甲的民办幼儿园符合财政补助的条件，依规定可申请10万元补助，但竞争很激烈。甲向民政局申请该补助。局长乙称：“申请的人很多，我优先考虑你，可以给你批10万元，但你拿到钱后分我4万元。”甲答应。甲领到10万元，然后送给乙4万元。乙构成受贿罪，数额是4万元。甲构成行贿罪，虽然甲符合条件，但谋取不正当竞争优势也属于谋取不正当利益。\CJKunderwave{注意}：就4万元而言，乙不构成贪污罪。因为甲符合补助条件，10万元是甲的合法所得。乙将10万元批给甲是符合规定的，乙不是非法将公款转移给甲及自已占有。

例2，甲的幼儿园符合财政补助的条件，依规定可申请10万元补助。甲向民政局申请该补助。局长乙称：“你可以申请14万元，我批给你，但你拿到钱后分我4万元。”甲答应。甲领到14万元后，将4万元送给乙。乙构成贪污罪，数额是4万元。甲构成贪污罪的帮助犯。\CJKunderwave{注意}：就4万元而言，表面上看甲是行贿，乙是受贿。但是，这4万元是乙非法将公款转移给甲，让甲先代管。甲只是帮助乙过了一下手，为乙的贪污提供了帮助。乙属于变相贪污。

% Source PDF physical page 415; see audit/source_coverage.json.
例3，甲的幼儿园不符合财政补助的条件。甲带着非法占有目的，向民政局申请10万元补助，为办成此事，送局长乙4万元。乙明知甲的幼儿园不符合条件，仍给甲批了10万元补助。甲构成行贿罪，乙构成受贿罪。乙还构成贪污罪，数额是10万元。贪污罪的非法占有目的既包括为自己非法占有，也包括为第三人非法占有。相应地，贪污行为既包括将公共财物转移给自己占有，也包括转移给第三人占有。甲构成贪污罪的共犯。最终，对甲以行贿罪和贪污罪（共犯）并罚，对乙以受贿罪和贪污罪（实行犯）并罚。

2.二者的竞合

例如，副乡长甲想升职，有求于县长乙，乙要求甲从乡财政中拿出30万元给自己。甲照办，事后也没归还30万元。甲将公款送人，构成贪污罪；送的行为构成行贿罪，想象竞合，择一重罪论处。乙构成贪污罪的教唆犯和受贿罪，想象竞合，择一重罪论处。

\subsubsection*{（三）受贿罪与敲诈勒索罪的关系}

[结论］二者不是对立排斥关系，而是中立关系，可以想象竞合。

例1，甲的公司有违法行为。官员乙威胁甲，不给30万元，就立案调查。甲很害怕，被迫给付。乙没有立案调查。乙构成受贿罪和敲诈勒索罪，想象竞合，择一重罪论处。\footnote{甲构成行贿罪。理由是，第一，虽被敲诈勒索，但是甲还保有选择自由，没有达到无法反抗的程度。第二，甲不属于受强制的紧急避险，因为紧急避险要求为保护较大合法权益，不得已损害较小的合法权益。甲想逃避立案调查，不属于合法权益。第三，甲实际获得了不正当利益。若甲没有实际获得不正当利益，则不构成行贿罪。第四，甲一方面是敲诈勒索罪的被害人，另一方面构成行贿罪。这并不矛盾。在程序上，基于甲是被害人，先将30万元退还给甲，然后由于是贿赂款，再没收掉。或者直接没收即可。}

例2，甲欲办理某项不正当业务，有求于官员乙。乙索要30万元，否则不予办理。甲予以给付。乙构成受贿罪，不构成敲诈勒索罪，因为索要行为不会使甲产生恐惧心理，不属于恐吓行为。甲构成行贿罪。

\subsubsection*{（四）受贿罪与诈骗罪的关系}

[结论］二者不是对立排斥关系，而是中立关系，可以想象竞合。

[提示］受贿罪中的许诺既可以是真心实意，也可以是虚假承诺，也即并不打算帮请托人办事。

例1，老板乙向官员甲送钱，甲收下钱后，乙便提出要办的违规事情，甲表示会给乙办事。实际上，甲是虚与委蛇，敷衍了事，并不想真心实意办事。甲构成受贿罪。但是，甲不构成诈骗罪，因为甲是在乙给钱后作出虚假承诺。乙构成行贿罪。（2017年第62题）

例2，老板乙求官员甲办违规事情，甲表示，要办事，需要给钱。乙给了钱。实际上，甲是虚假许诺，并不想给办事。甲构成诈骗罪，因为甲带着非法占有乙财物的目的，对乙实施了欺骗，使乙基于认识错误而处分财物。不过，不能因为甲构成诈骗罪，就认为甲不构成受贿罪。甲能欺骗成功，是因为甲基于职务便利。甲亵渎了职务，所以甲同时构成受贿罪。甲的一个行为同时触犯两罪，想象竞合，择一重罪论处。\footnote{多数说认为，乙不构成行贿罪。理由是，行贿罪的保护法益是国家工作人员职务行为的不可收买性。由于甲一开始就不想办事，只想欺骗乙，所以乙的行为不具有收买甲的职务行为的可能性。乙是诈骗罪的被害人。}

% Source PDF physical page 416; see audit/source_coverage.json.
\section{行贿罪}

[问题］王某符合晋升条件，但是为了使自已优于他人，向组织部部长李某送5万元。李某收受，但是因为其他领导反对，未能将事办成。王某是否构成行贿罪？（跑官买官案）（2005年第65题）

第389条为谋取\CJKunderwave{不正当利益}，给予国家工作人员以财物的，是行贿罪。

在经济往来中，违反国家规定，给予国家工作人员以财物，数额较大的，或者违反国家规定，给予国家工作人员以各种名义的回扣、手续费的，以行贿论处。

\CJKunderwave{因被勒索}给予国家工作人员以财物，\CJKunderwave{没有获得不正当利益的}，不是行贿。

第390条对犯行贿罪的，\CJKunderwave{处三年以下有期徒刑（旧规定是五年以下有期徒刑）}或者拘役，并处罚金；因行贿谋取不正当利益，情节严重的，或者使国家利益遭受重大损失的，处三年以上十年以下有期徒刑，并处罚金；情节特别严重的，或者使国家利益遭受特别重大损失的，处十年以上有期徒刑或者无期徒刑，并处罚金或者没收财产。

有下列情形之一的，从重处罚：·.··

行贿人在被追诉前主动交待行贿行为的，可以从轻或者减轻处罚。其中，犯罪较轻的，对调查突破、侦破重大案件起关键作用的，或者有重大立功表现的，可以减轻或者免除处罚。[本条经《刑法修正案（十二）》修订]

\subsection*{一、构成要件}

\subsubsection*{（一）“为谋取不正当利益”}

行贿罪的行为公式：为谋取不正当利益，给予国家工作人员以财物。

1.“为谋取不正当利益”的要求

主观上，要求具有谋取不正当利益的目的。

客观上，要求利益具有不正当性。例如（认识错误案），政府应当给甲的企业补偿50万元，甲以为其企业不符合补偿条件，带着谋取“不正当利益”的目的，希望主管官员乙给补偿，给乙送钱。乙收钱。乙构成受贿罪。甲不构成行贿罪，因为甲谋取的利益在客观上属于正当利益。

2.“为谋取不正当利益”的范围

第一，不正当利益不限于非法利益，获取不公平的竞争优势也属于不正当利益。例如，本节开头的“跑官买官案”中，王某属于谋取不正当利益，构成行贿罪。李某构成受贿罪。

第二，公民请求国家工作人员依法办事，属于正当利益。

\subsubsection*{（二）行为类型}

1.送钱的缘起

（1）主动给钱。这种情形下，成立行贿罪，要求带有“为谋取不正当利益”的目的，但不要求事后实际获得不正当利益。即使没有实际获得不正当利益，也构成行贿罪，而且只要钱给出去，行贿罪便既遂。

（2）被动给钱，也即因被勒索而给钱。这种情形下，成立行贿罪，不仅要求产生“为谋取不正当利益”的目的，而且要求实际获得了不正当利益，才构成行贿罪。法条依

% Source PDF physical page 417; see audit/source_coverage.json.
据是本罪第3款规定：“因被勒索给予国家工作人员以财物，没有获得不正当利益的，不是行贿。”该规定从反面否定犯罪成立，属于消极的构成要件要素（2010年第65题）。该规定因为有利于被告人，因此可以类推适用于对非国家工作人员行贿罪、对单位行贿罪。

2.送钱的时间

（1）事前有约定，事前或事后送钱。例如，税务局局长甲与老板乙达成约定，甲“办事”，乙送钱。可以先送钱，后办事，也可以先办事，后送钱。甲构成受贿罪，乙构成行贿罪。可见，只要事前有约定，权钱交易就订立，受贿罪、行贿罪就成立，至于何时送钱、收钱，不影响犯罪成立，只影响犯罪既遂。

（2）事前无约定，事后送钱。例如，税务局局长甲没有受贿意图，为乙的公司办理了不正当的业务。事后，乙欲感谢甲，甲说：“我办事时没想图你们的钱，就是觉得你们不容易。”乙说：“正因如此，更要感谢您！”便送甲4万元，甲予以接受。甲成立受贿罪，属于事后受贿。乙是否构成行贿罪？构成。理由是，行贿罪的“为谋取不正当利益而送钱”既包括为谋取未来的不正当利益而送钱，也包括为了已经谋取的不正当利益而送钱。乙客观上谋取到了不正当利益，乙为此送钱，甲乙形成权钱交易。\footnote{少数说认为，不构成。理由是，行贿罪的“为谋取不正当利益而送钱”，仅限于为谋取未来的不正当利益。}

\subsubsection*{（三）既遂问题}

既遂标准：国家工作人员客观上占有了财物。

例1，甲送官员乙3万元，乙当场拒绝。甲成立行贿罪未遂，乙不构成受贿罪。

例2，甲将一张存有3万元的银行卡悄悄放在官员乙家茶几底下。乙次日发现，便立即上交。甲构成行贿罪既遂，乙不构成受贿罪。

[注意］主动给钱型行贿罪，要构成既遂，不要求实际获得了不正当利益。例如，本节开头的“跑官买官案”中，虽然王某最终没能升官，但是仍构成行贿罪既遂。

\subsubsection*{（四）数额问题}

行贿罪是故意犯罪，成立条件：行贿人带着行贿数额较大财物（3万元）的故意，实施了行贿行为。既遂条件：将数额较大财物送给国家工作人员，由其客观上占有。\footnote{2016年4月18日最高人民法院、最高人民检察院《关于办理贪污贿赂刑事案件适用法律若干问题的解释》规定，行贿罪的“数额较大”是3万元。之前的司法解释规定，行贿罪的“数额较大”是1万元。}

例如，甲花了50万元买了一幅名画，送给官员乙，声称价值50万元。实际甲被卖家骗了，该画是品，只值1万元。甲具有行贿数额较大财物的故意和行为，构成行贿罪，但是客观上没有行贿数额较大的财物，因此构成行贿罪未遂。乙构成受贿罪未遂。

[注意］行贿数额=受贿数额。例如，局长甲为乙办过事。甲装修房子，装修公司报价100万元，乙对甲说：“你只需出40万元，剩余的你不用管。”甲答应。乙对装修公司老板丙讲：“我给你30万元，我以后给你介绍个大工程，剩余30万元能不能免掉？”丙答应。甲获得的财产性利益是60万元，受贿数额是60万元。基于此，乙的行贿数额也应是60万元。在此不能以乙实际支出的30万元来计算，不能认为“乙的行贿数额是30万元”。\footnote{有时候，行贿方不构成行贿罪，或者受贿方不构成受贿罪。此时，等号原理仅就实际数额而言。例如，官员勒索老板，老板被迫给了10万元，没有获得不正当利益。老板不构成行贿罪，但老板支付的10万元等于官员受贿的10万元。}

% Source PDF physical page 418; see audit/source_coverage.json.
\subsection*{二、认定问题}

\subsubsection*{（一）行贿罪与受贿罪的关系}

行贿罪与受贿罪属于对向犯。一般情况下，双方同时成立犯罪。例外情况下，只有一方成立犯罪。具体情形有：

1.因被勒索给予财物，没有获得不正当利益的，不是行贿。但国家工作人员的行为仍然是索取贿赂，成立受贿罪。

2.为了谋取正当利益而给予国家工作人员以财物的，不是行贿；但国家工作人员接受财物的行为成立受贿罪。

3.为了谋取不正当利益而给予国家工作人员以财物的，构成行贿罪；但国家工作人员拒绝的，不构成受贿罪。

\subsubsection*{（二）罪数问题}

“谋取不正当利益”若构成犯罪，则存在罪数问题。

1.原则上，行贿人谋取的不正当利益构成犯罪，应与行贿罪并罚。

例1，甲送给国有电力公司经理乙3万元，让乙为其盗窃电力，乙利用职权为其照办。甲盗窃了电力，价值10万元。甲构成行贿罪和盗窃罪，并罚。乙构成受贿罪和贪污罪（为第三人非法占有电力），并罚。

例2，甲送给税务局局长乙3万元，让乙为自已逃税了50万元。乙照办。甲构成行贿罪和逃税罪，并罚。乙构成受贿罪和徇私舞弊不征、少征税款罪，并罚。

2.如果出现重复处罚，应坚持禁止重复处罚原则，不得数罪并罚。

例如，甲送给警察乙3万元，让乙为自已办理一个假户口。甲构成行贿罪，不再定买卖国家机关证件罪，因为甲给乙3万元，就相当于从乙处买了个假户口，若再定买卖国家机关证件罪，就重复处罚了。

\section{斡旋受贿}

第388条国家工作人员\CJKunderwave{利用本人职权或者地位形成的便利条件，通过其他国家工作}\CJKunderwave{人员职务上的行为}，为请托人谋取\CJKunderwave{不正当利益}，索取请托人财物或者收受请托人财物的，以受贿论处。

斡旋受贿，是指国家工作人员将自己的翰旋行为与他人的财物结成不正当的对价关系。

\InsertDiagram{figures/manual_p418_1.tex}

\subsection*{一、基本要点}

1.行为主体（翰旋者）：国家工作人员，不包括单位。

2.罪名不是斡旋受贿罪，而是受贿罪。

3.为请托人谋取的是不正当利益。这表明，本罪的保护法益有两个：直接保护的法益是翰旋者职权影响力的不可收买性，间接保护的法益是终端办事人职务行为的公正性。假如不要求谋取不正当利益，则保护法益就没有“终端办事人职务行为的公正性”。

% Source PDF physical page 419; see audit/source_coverage.json.
\subsection*{二、构成要件}

犯罪成立条件：许诺为请托人实施翰旋行为、谋取不正当利益，基于此索取或收受财物。犯罪既遂条件：收到数额较大的财物。（2023年试题，2026年试题）

[注意］只要求许诺实施斡旋行为，不要求实际实施了翰旋行为，不要求终端办事人许诺办事，不要求终端办事人办成事。这是因为，请托人花钱买的是翰旋者的翰旋行为，而非终端办事人的办事行为。只要斡旋者许诺“办事”，就会让请托人认识到，翰旋者的职权影响力是可以花钱收买的，那么，保护法益“翰旋者职权影响力的不可收买性”就遭到了破坏。（2026年司法解释赞成该观点）

例如，甲请财政局局长乙向税务局局长丙说情，为自已公司非法减免税款。乙许诺，并收了甲的钱。之后，乙并没有向丙说情，或者说情了，丙拒绝了，或者丙答应了，但没办成事。这些情形下，乙均构成受贿罪既遂。甲构成行贿罪既遂。\footnote{注意：“为请托人谋取不正当利益”，只是一种主观要求，不要求客观上存在不正当利益。

例如，甲的表弟乙涉嫌犯罪被捕，甲请法官丙帮忙“违规捞人”，事成后给丙20万元。丙许诺“给钱，我就办事”。丙找到主审法官丁，打听案情。丁说：“根据法律规定，乙无罪。”丙便欺骗甲：“这事我办成了，过几天就放人，到时给钱啊。”几日后，乙被宣告无罪。甲给丙20万元。

第一，丙许诺“给钱就办事”，就成立斡旋受贿。虽然客观上利益是正当的，但丙仍成立斡旋受贿。这表明，成立斡旋受贿，间接保护的法益“终端办事人职务行为的公正性”，不要求遭受现实侵害、具体危险，只要求遭受抽象危险。当丙许诺“给钱，就为你谋取不正当利益”，该间接法益就遭受了抽象危险。

第二，在这一点上，斡旋受贿与行贿罪有所不同。行贿罪中的“为谋取不正当利益”既要求是主观目的，也要求客观存在不正当利益。由于乙是无罪的，所以甲谋取的利益在客观上是正当的。因此，甲不构成行贿罪。

第三，丙对甲还构成诈骗罪，与斡旋受贿想象竞合，择一重罪论处。甲是诈骗罪的受害人。}

\subsection*{三、共犯问题}

以下模型是最完整的行为模型，比较复杂，但很重要，希望勇敢的你耐心学习，多看两遍。

共同前提事实：甲为了让税务局局长丙非法减免税款，送给本县财政局局长乙100万元，对乙讲：“这笔钱，你自己看着办，只要办成事就行。”乙答应办事，收了100万元。

1.模型一：乙向丙游说，给丙60万元。丙收了钱，许诺办事。

（1）丙构成普通的受贿罪，既遂数额是60万元。

（2）乙向丙送钱这个行为，第一，构成丙的受贿罪的共犯；第二，构成行贿罪，与甲构成行贿罪的共同正犯。两罪想象竞合，择一重罪论处。

（3）乙收甲的钱，向甲承诺办事，构成斡旋受贿型受贿罪。既遂数额不是100万元，而是100万元-60万元=40万元。这是因为，乙收下100万元时，这笔钱的去向和分配还是待定状态。当乙给丙60万元，表明乙只想收甲40万元。假如乙给了丙100万元，则乙的斡旋受贿型受贿罪就没有既遂数额，构成犯罪未遂。

（4）乙的罪数问题。乙收甲钱的行为与乙给丙送钱的行为，这两个行为属于原因与结果的牵连犯关系，择一重罪论处。\footnote{参见张明楷编著：《刑法的私垫（之四）》，北京大学出版社2025年版，第287页。乙收甲钱的行为和乙向丙送钱的行为是两个独立行为，前者没有涵盖后者。这是因为，乙给丙送钱的行为不属于翰旋行为（说情请托行为）的必备内容，因为实施翰旋行为并不要求送钱，送钱是翰旋行为之外的行为。例如，乙可以给丙说：“甲是我的表弟，请照顾”，没给丙送钱。丙考虑到甲是乙的表弟，没收钱，给办了事。}

% Source PDF physical page 420; see audit/source_coverage.json.
（5）甲送钱时，对钱的去向（乙独吞还是送给丙）没有明确指示，表明甲既有向乙行贿的意图，也有向丙行贿的概括故意。第一，甲对乙构成行贿罪，既遂数额是40万元。这是因为，受贿数额=行贿数额。两个数额要能划等号。第二，甲对丙构成行贿罪，甲乙是这个行贿罪的共同正犯。既遂数额是60万元。甲的两个送钱行为（一个是甲给乙送钱，一个是乙代表甲给丙送钱）构成两个行贿罪，属于连续犯，定一个行贿罪，数额累计为100万元。

2.模型二：乙向丙游说，对丙说自己收了甲100万元，希望丙不要收这笔钱，但希望丙办事，丙知道了，没向乙要钱，但仍办事。

（1）丙构成受贿罪，既遂数额是100万元。这是因为，\CJKunderwave{受贿罪的收钱，既包括为自己}\CJKunderwave{收钱，也包括为第三人收钱，亦即要对第三人收的钱负责}。例如，丙的情妇告知丙：“我收了甲的100万元，我不给你分钱，你能不能办事？”丙答应并办事。丙构成受贿罪，既遂数额是100万元。这里的既遂数额100万元，不是指自然事实角度丙收到了100万元，而是指在法律评价角度，丙要对情妇收的100万元负责。

（2）乙针对丙，构成丙的受贿罪的共犯，既遂数额也是100万元。乙针对丙，不构成行贿罪，因为乙没有送钱给丙的意图。

（3）乙针对甲，构成斡旋受贿型受贿罪，既遂数额是100万元。\footnote{可能有人认为，由于100万元算是丙的受贿罪的既遂数额，所以这100万元就不能算到乙的头上，乙的受贿罪就没有既遂数额，只能视为未遂。然而，所谓丙的受贿数额是100万元，是指丙要对乙收的100万元负责，并非指丙在事实上收了100万元。而乙在事实上收了100万元，所以乙的受贿罪既遂数额是100万元。假如，乙将100万元给了丙，丙在事实上收了100万元，则乙的受贿罪数额只能是零，构成受贿罪未遂。概言之，在自然事实层面，乙收到的钱与送出的钱只能是加减关系，但在法律评价角度，同一笔钱的数额不一定是加减关系，可以重复评价，只要不重复处罚即可。}

（4）乙的罪数问题。“乙收甲的钱”这个行为构成斡旋受贿型受贿罪，“\CJKunderwave{乙告知丙自}\CJKunderwave{已收了}甲的钱”，这个告知行为构成丙的受贿罪的共犯。这两个行为属于原因与结果的牵连犯关系，择一重罪论处。

（5）甲送钱时，对乙和丙都有行贿的概括故意。第一，甲对丙构成行贿罪，既遂数额是100万元。第二，甲对乙构成行贿罪，既遂数额是100万元。甲的两个行贿罪是想象竞合，择一重罪论处。这是因为，甲只有一个送钱行为，也即甲给乙送钱，乙没有代表甲给丙送钱。\footnote{可能有人认为，甲只送出去一份100万元，为何两个行贿罪的既遂数额都是100万元？是不是对100万元重复评价了？是重复评价了。这是因为，丙和乙的受贿既遂数额都是100万元，这就是从法律评价角度重复评价了，由此导致对行贿方的行贿罪数额也就重复评价了。想象竞合犯就是对一个行为进行重复评价，但没有重复处罚，而是择一重罪论处。所谓“禁止重复评价原则”，实际上是指禁止重复处罚。}

3.模型三：乙向丙游说，隐瞒自己收了钱，但丙猜出来乙收了甲100万元，仍办事。（2023年主观题）

（1）丙构成受贿罪，既遂数额是100万元。这是法律评价角度的既遂数额，意思是对乙收的100万元负责。

（2）乙针对丙，不构成受贿罪的共犯，因为乙没有帮助丙受贿的故意。乙针对丙，也不构成行贿罪，因为乙没有向丙行贿的意图。

（3）乙针对甲，构成斡旋受贿型受贿罪，既遂数额是100万元。

% Source PDF physical page 421; see audit/source_coverage.json.
（4）甲送钱时，对乙和丙都有行贿的概括故意。第一，甲对丙构成行贿罪，既遂数额是100万元。第二，甲对乙构成行贿罪，既遂数额是100万元。甲的两个行贿罪是想象竞合，择一重罪论处。假如甲给乙送钱时说：“这100万元就是专门给你的，如果终端办事人要钱，你再跟我说”，那么，甲只对乙构成行贿罪，对丙不构成行贿罪。（2023年主观题）\footnote{可能有人认为，此时丙构成受贿罪，但没人对丙构成行贿罪，是否不合理？答复是，丙构成受贿罪，是指要对乙的受贿结果负责，只要乙有对应的行贿人就可以了。}

\subsection*{四、斡旋受贿与普通受贿的区别}

如果斡旋者是办事人的上级，具有隶属、制约关系，则不构成斡旋受贿，而属于普通受贿。这是因为，斡旋行为是指游说，请求给个面子。而上级给下级发话，就不是游说，而带有命令要求意味。这种隶属、制约关系不限于主管关系，不限于直接的上下级关系。

例如，甲为了让税务局局长非法减免税款，送给本县主管税务的副县长或主管公安的副县长乙一笔钱，让办此事。乙答应，向税务局局长丙告知此事，丙答应照办。乙属于普通受贿，构成受贿罪。注意：上级领导既包括直接主管者，也包括非直接主管者，但必须是同一辖区的领导。如果乙是隔壁县的副县长，则不属于普通受贿，而属于斡旋受贿。

[总结】上级找下级，是普通受贿；同级找同级，是斡旋受贿；下级找上级，是斡旋受贿。

\section{利用影响力受贿罪}

第388条之一国家工作人员的\CJKunderwave{近亲属}或者其他与该国家工作人员\CJKunderwave{关系密切的人}，\CJKunderwave{通过}该国家工作人员职务上的行为，或者利用该国家工作人员职权或者地位形成的便利条件，\CJKunderwave{通过}其他国家工作人员职务上的行为，为请托人谋取\CJKunderwave{不正当利益，}索取请托人财物或者收受请托人财物，数额较大或者有其他较重情节的，处三年以下有期徒刑或者拘役，并处罚金；….

\CJKunderwave{离职}的国家工作人员或者其\CJKunderwave{近亲属}以及其他与其\CJKunderwave{关系密切的人}，利用该离职的国家工作人员\CJKunderwave{原职权}或者地位形成的便利条件实施前款行为的，依照前款的规定定罪处罚。

\InsertDiagram{figures/manual_p421_1.tex}

利用影响力受贿罪与对有影响力的人行贿罪，是对向犯关系。重点讲解利用影响力受贿罪。

\subsection*{一、行为主体}

行为主体有三种：（1）国家工作人员的近亲属或关系密切人。（2）离职的国家工作人员。（3）离职的国家工作人员的近亲属或关系密切人。

[注意］关系密切人，是指与国家工作人员或离职的国家工作人员具有利益关系或利害关系的人。其中的利益关系不限于物质利益，还可以包括其他利益，例如，情人关系、前妻关系、秘书、司机等。其中的利害关系包括制约关系，如握有国家工作人员把柄，依此来制约国家工作人员。也即“关系密切人”不等于“关系好的人”。而斡旋受贿的中间人不能对终端办事人有制约关系，否则就是普通受贿。

% Source PDF physical page 422; see audit/source_coverage.json.
\subsection*{二、构成要件}

犯罪成立条件：行为人为请托人实施游说行为，并且，终端办事人许诺为请托人\CJKunderwave{谋取}\CJKunderwave{不正当利益}，行为人基于此索取或收受财物。犯罪既遂条件：收到数额较大的财物。（2026年试题）

这是因为，本罪的保护法益不是中间人的职权影响力，因为中间人不必是国家工作人员，不要求有职权。本罪的保护法益是\CJKunderwave{终端办事人职务行为的不可收买性和公正性}。有“公正性”，是因为本罪要求为请托人谋取不正当利益。

例如，甲为了让税务局局长丙非法减免税款，送给税务局局长的司机乙一笔钱，让办此事。乙答应，向税务局局长丙游说，丙许诺办事。乙构成利用影响力受贿罪。注意：如果乙未实施游说行为，或者实施了游说，但丙拒绝办事，则乙不构成利用影响力受贿罪。\CJKunderwave{这一点，利用影响力受贿罪与斡旋受贿是不同的}。（2020年、2023年试题）\footnote{注意：“为请托人谋取不正当利益”，不仅是一种主观要求，而且要求，在终端办事人许诺时，客观上可能存在不正当利益。这一点，本罪与斡旋受贿也是不同的。

例如，甲的表弟乙涉嫌犯罪被捕，甲请主审法官丁的情妇丙帮忙“违规捞人”，事成后给丙20万元。丙答应办事。丙向丁游说。

第一种情形，丁说：“案子还在审，我会办成事。”此时，丙就成立利用影响力受贿罪。过段时间，丁说：“案子审清楚了，按照法律，乙无罪，不需要我特别帮忙。”此时，不影响丙成立本罪的结论。然后，丙欺骗甲：“我办成事了，过几天就宣判乙无罪，给钱吧！”甲不知情而给钱。丙构成本罪既遂和诈骗罪既遂，想象竞合。

第二种情形，丁说：“案子审完了，按照法律，乙无罪。”由于不存在“丁许诺为甲谋取不正当利益”，所以，丙不构成利用影响力受贿罪。不过，丙又欺骗甲：“这事我办成了，过几天就放人，给钱吧。”甲不知情而给钱。丙构成诈骗罪既遂。

[总结]谋取不正当利益的要求。

（1）行贿罪：主观要求有该目的；客观要求存在不正当利益。

（2）斡旋受贿：主观要求有该目的；客观不要求存在不正当利益。

（3）利用影响力受贿罪：主观要求有该目的；客观要求，在终端办事人许诺时，可能存在不正当利益。}

\subsection*{三、共犯问题}

共同前提事实：甲为了让税务局局长丙非法减免税款，送给丙的司机乙100万元，“这笔钱，你自己看着办，只要办成事就行”。乙答应办事，收了100万元。

1.模型一：乙向丙游说，给丙60万元。丙收了钱，许诺办事。

（1）丙构成受贿罪，既遂数额是60万元。

（2）乙向丙送钱这个行为，第一，构成丙的受贿罪的共犯；第二，构成行贿罪，与甲构成行贿罪的共同正犯。两罪想象竞合，择一重罪论处。

（3）当丙承诺办事时，乙针对甲，成立利用影响力受贿罪。既遂数额不是100万元，而是100万元-60万元=40万元。这是因为，乙收下100万元时，这笔钱的去向和分配还是待定状态。当乙给丙60万元，表明乙只想收甲40万元。假如乙给了丙100万元，则乙的利用影响力受贿罪就没有既遂数额，构成犯罪未遂。

（4）乙的罪数问题。乙收甲钱的行为与乙向丙送钱的行为，这两个行为属于原因与结果的牵连犯关系，择一重罪论处。\footnote{乙收甲钱的行为和乙向丙送钱的行为是两个独立行为，前者没有涵盖后者。这是因为，第一，乙向丙送钱的行为不是乙向丙实施游说行为（说情请托行为）的必备内容，实施游说行为并不要求送钱，送钱是游说行为之外的行为。

例如，乙可以给丙说：“甲是我的表弟，请照顾”，没给丙送钱。丙考虑到甲是乙的表弟，没收钱，给办了事。第二，当乙向丙请托，丙许诺办事时，乙收甲钱的行为才成立利用影响力受贿罪。但是，这不意味着，只有乙向丙送钱、丙收钱时，乙才成立利用影响力受贿罪。所以，“乙收甲钱成立利用影响力受贿罪”与“乙向丙送钱”是两个独立行为。}

% Source PDF physical page 423; see audit/source_coverage.json.
（5）甲送钱时，对钱的去向（乙独吞还是送给丙）没有明确指示，表明甲既有向乙行贿的意图，也有向丙行贿的概括故意。第一，甲对乙构成对有影响力的人行贿罪，既遂数额是40万元。这是因为，受贿数额=行贿数额。两个数额要能划等号。第二，甲对丙构成行贿罪，甲乙是这个行贿罪的共同正犯。既遂数额是60万元。甲的两个送钱行为构成两个行贿犯罪，罪名不同，不是连续犯，应两罪并罚。

2.模型二：乙向丙游说，对丙说自己收了甲100万元，希望丙不要收这笔钱，但希望丙办事，丙知道了，没向乙要钱，但仍办事。

（1）丙构成受贿罪，既遂数额是100万元。这是因为，\CJKunderwave{受贿罪的收钱，既包括为自己}\CJKunderwave{收钱，也包括为第三人收钱}。认定丙的受贿罪既遂数额是100万元，不是指丙在事实上收了100万元，而是指丙要对乙收的100万元负责。这是一种法律评价角度的既遂数额。

（2）乙针对丙，构成丙的受贿罪的共犯，既遂数额也是100万元。乙针对丙，不构成行贿罪，因为乙没有送钱给丙的意图。

（3）乙针对甲，构成利用影响力受贿罪，既遂数额是100万元。

（4）乙的罪数问题。“乙收甲的钱”这个行为构成利用影响力受贿罪，“乙告知丙自已收了甲的钱”，这个告知行为构成丙的受贿罪的共犯。这两个行为属于原因与结果的牵连犯关系，择一重罪论处。

（5）甲送钱时，对乙和丙都有行贿的概括故意。第一，甲对丙构成行贿罪，既遂数额是100万元。第二，甲对乙构成对有影响力的人行贿罪，既遂数额是100万元。甲的两个行贿犯罪是想象竞合，择一重罪论处。

3.模型三：乙向丙游说，隐瞒自已收了钱，但丙猜出来乙收了甲100万元，仍办事。

（1）丙构成受贿罪，既遂数额是100万元。这是法律评价角度的既遂数额，意思是对乙收的100万元负责。

（2）乙针对丙，不构成受贿罪的共犯，因为乙没有帮助丙受贿的故意。乙针对丙，也不构成行贿罪，因为乙没有向丙行贿的意图。

（3）乙针对甲，构成利用影响力受贿罪，既遂数额是100万元。

（4）甲送钱时，对乙和丙都有行贿的概括故意。第一，甲对丙构成行贿罪，既遂数额是100万元。第二，甲对乙构成对有影响力的人行贿罪，既遂数额是100万元。甲的两个行贿犯罪是想象竞合，择一重罪论处。

\subsection*{四、总结：两个中间人}

模型1：请托人（甲）\ensuremath{\rightarrow}中间人（乙，国家工作人员）\ensuremath{\rightarrow}中间人（丙，国家工作人员）\ensuremath{\rightarrow}终端办事人（丁）。

例如，甲欲让税务局局长丁违规办事，送给乙（财政局局长）100万元，告知：“这么多钱，这么个事，你看着办。”乙收钱，送给丙（市场监督局局长）70万元，让帮忙找丁，告知“这么多钱，这么个事，你看着办”。丙收钱，送给丁40万元，丁收钱并办事。

（1）丁构成受贿罪（40万元）。（2）丙构成斡旋受贿型受贿罪（30万元），构成丁的受贿罪的共犯，构成行贿罪的共同正犯。（3）乙构成斡旋受贿型受贿罪（30万元），构成丙的受贿罪的共犯，构成丁的受贿罪的共犯，构成行贿罪的共同正犯。（4）甲对乙构成行贿罪（30万元），对丙构成行贿罪（30万元），对丁构成行贿罪（40万元）。

模型2：请托人（甲）\ensuremath{\rightarrow}中间人（乙，非国家工作人员）\ensuremath{\rightarrow}中间人（丙，国家工作人员）\ensuremath{\rightarrow}终端办事人（丁）。

例如，甲欲让税务局局长丁违规办事，送给乙（财政局局长丙的司机）100万元，告

% Source PDF physical page 424; see audit/source_coverage.json.
知：“这么多钱，这么个事，你看着办。”乙收钱，送给丙70万元，让帮忙找丁，告知“这么多钱，这么个事，你看着办”。丙收钱，送给丁40万元，丁收钱并办事。（1）丁构成受贿罪（40万元）。（2）丙构成斡旋受贿型受贿罪（30万元），构成丁的受贿罪的共犯，构成行贿罪的共同正犯。（3）乙构成利用影响力受贿罪（30万元），构成丙的受贿罪的共犯，构成丁的受贿罪的共犯，构成行贿罪的共同正犯。（4）甲对乙构成对有影响力的人行贿罪（30万元），对丙构成行贿罪（30万元），对丁构成行贿罪（40万元）。（2009年第64题，2023年试题）

模型3：请托人（甲）\ensuremath{\rightarrow}中间人（乙，国家工作人员）\ensuremath{\rightarrow}中间人（丙，非国家工作人员）\ensuremath{\rightarrow}终端办事人（丁）。

例如，甲欲让税务局局长丁违规办事，送给乙（财政局局长）100万元，告知：“这么多钱，这么个事，你看着办。”乙收钱，送给丙（丁的司机）70万元，让帮忙找丁，告知“这么多钱，这么个事，你看着办”。丙收钱，送给丁40万元，丁收钱并办事。（1)丁构成受贿罪（40万元）。（2）丙构成利用影响力受贿罪（30万元），构成丁的受贿罪的共犯，构成行贿罪的共同正犯。（3）乙构成斡旋受贿型受贿罪（30万元），构成丙的利用影响力受贿罪的共犯，构成丁的受贿罪的共犯，构成对有影响力的人行贿罪的共同正犯，构成行贿罪的共同正犯。（4）甲对乙构成行贿罪（30万元），对丙构成对有影响力的人行贿罪（30万元），对丁构成行贿罪（40万元）。

\subsection*{五、利用影响力受贿罪与斡旋受贿的区别}

二者的相同点是：都是中间人受贿，都为他人谋取不正当利益。二者的区别：

1.主体身份不同。前者，不是国家工作人员。后者，是国家工作人员。

2.对价关系不同。前者是终端办事人（国家工作人员）的职务行为与请托人钱财的对价关系。后者是翰旋者（国家工作人员）的职务地位与请托人钱财的对价关系。

3.发挥影响力的根据不同（关键区别）。二者的行为主体都是中间人，都对终端办事人发挥影响力，但发挥影响力的根据不同。前者的根据是私人关系；后者的根据是国家工作人员的职权地位及工作关系。如果两种影响力均有，则二者想象竞合，择一重罪论处。

例如，甲为了让税务局局长非法减免税款，送给本县财政局局长乙（税务局局长的情妇）一笔钱，让办此事。乙答应，并找税务局局长丙游说，丙答应。乙同时触犯利用影响力受贿罪和斡旋受贿，想象竞合，择一重罪论处。

[总结1]行贿与受贿的罪名关系。（2009年第20题）

\InsertDiagram{figures/manual_p424_1.tex}

% Source PDF physical page 425; see audit/source_coverage.json.
[总结2]谋取利益的正当与不正当问题。

1.所有的行贿罪，要求“为谋取不正当利益”；包括五个罪名：行贿罪，对有影响力的人行贿罪，单位行贿罪，对单位行贿罪，对非国家工作人员行贿罪。

2.受贿罪中：（1）普通的受贿罪，单位受贿罪，非国家工作人员受贿罪，这三个犯罪中，“为他人谋取利益”包括正当利益和不正当利益。（2）斡旋受贿型受贿罪，利用影响力受贿罪，这两个犯罪中，要求“为他人谋取不正当利益”。

[结论】除了三个罪（普通的受贿罪，单位受贿罪，非国家工作人员受贿罪），其他罪都要求谋取不正当利益。\footnote{单位受贿罪、对单位行贿罪的“单位”要求是国有单位。单位行贿罪的“单位”，没有此要求。}

以下为普通罪名

1.介绍贿赂罪（第392条）。这是指，在请托人和国家工作人员之间沟通关系、合条件，使贿赂行为得以实现的行为。

例如，政治客甲在老板乙和官员丙之间牵线搭桥，促成乙丙之间的权钱交易。甲构成介绍贿赂罪，也构成乙的行贿罪的共犯、丙的受贿罪的共犯，三罪想象竞合，择一重罪论处。（2026年司法解释的观点）

2.巨额财产来源不明罪

第395条第1款国家工作人员的财产、支出明显超过合法收入，差额巨大的，可以责令该国家工作人员说明来源，\CJKunderwave{不能说明来源的}，差额部分以非法所得论，处五年以下有期徒刑或者拘役；差额特别巨大的，处五年以上十年以下有期徒刑。财产的差额部分予以追缴。

巨额财产来源不明罪，是指国家工作人员的财产或者支出明显超过合法收入，差额巨大，不能说明其合法来源的行为。

（1）实行行为。“国家工作人员的财产、支出明显超过合法收人，差额巨大”不是本罪的实行行为，而是前提条件。本罪的实行行为是责令说明来源，不能说明来源。也即，本罪是真正不作为犯。

（2）认定问题：

①行为人拥有巨额财产，但不能说明合法来源的，应认定为本罪。

②行为人拥有巨额财产，只要在一审判决前，本人说明了其合法来源，就不构成本罪；如果说明了其非法来源，并查证属实的，就按其行为性质认定犯罪，不认定为本罪。（2012年第63题）

③行为人拥有巨额财产，本人不能说明其合法来源的，人民法院判决成立本罪；但司法机关后来查清了该巨额财产的来源：如果来源是合法的，原来的判决必须维持，不能更改；如果来源是非法的，则按非法来源的性质再次定罪，也不能推翻原来的判决。

% Source PDF physical page 426; see audit/source_coverage.json.
\OriginalPanel{典型真题}

关于受贿罪，下列哪些选项是正确的？（2017年·卷二·62题）\footnote{[答案]ABCD。}A.国家工作人员明知其近亲属利用自己的职务行为受贿的，构成受贿罪B.国家工作人员虚假承诺利用职务之便为他人谋利，收取他人财物的，构成受贿罪C.国家机关工作人员实施渎职犯罪并收受贿赂，同时构成渎职罪和受贿罪的，除《刑法》有特别规定外，以渎职罪和受贿罪数罪并罚D.国家工作人员明知他人有请托事项而收受其财物，视为具备“为他人谋取利益”的构成要件，是否已实际为他人谋取利益，不影响受贿的认定

